\setlength{\absparsep}{18pt}
\begin{resumo}[Abstract]
    Probabilistic models for football outcomes have been extensively studied for European competitions, but applications to Brazilian football remain limited.
    This dissertation introduces a comprehensive dataset of Brazilian football matches spanning 2013 to 2025, covering the four national divisions and the \textit{Copa do Brasil}, with detailed information on players, goals, cards, and substitutions.
    Based on this dataset, we develop a collection of Bayesian models from the Bradley-Terry and Poisson families, formulated at both club and player levels.
    All models are validated through Simulation-Based Calibration to ensure that the inference procedure correctly approximates the posterior distribution.
    The validated models are then applied to \textit{Brasileirão Série A} (the Brazilian first division) data, where their predictive performance is compared using proper scoring rules.
    Finally, we extend the best-performing club-level specifications to the player level, enabling a more granular analysis of team performance.
    All data and code are publicly available to support reproducibility.

    \vspace{\onelineskip}

    \noindent
    \textbf{Keywords}: Brazilian football, Bayesian inference, Bradley-Terry model, Poisson model, Simulation-Based Calibration, player-level modelling.    
\end{resumo}

\begin{resumo}[Resumo]
    \begin{otherlanguage*}{portuguese}
        Modelos probabilísticos para resultados de futebol têm sido amplamente estudados para competições europeias, porém aplicações ao futebol brasileiro permanecem limitadas.
        Este trabalho apresenta um conjunto de dados abrangente de partidas do futebol brasileiro entre 2013 e 2025, cobrindo as quatro divisões nacionais e a Copa do Brasil, com informações detalhadas sobre jogadores, gols, cartões e substituições.
        Com base nesse conjunto de dados, desenvolvemos uma coleção de modelos bayesianos das famílias Bradley-Terry e Poisson, formulados tanto em nível de clube quanto de jogador.
        Todos os modelos são validados por meio de Calibração Baseada em Simulação para garantir que o procedimento de inferência reproduza corretamente a distribuição a posteriori.
        Os modelos validados são então aplicados aos dados do Brasileirão Série A (a primeira divisão do futebol brasileiro), onde seu desempenho preditivo é comparado utilizando regras de pontuação próprias.
        Por fim, estendemos as especificações de nível de clube com melhor desempenho para o nível de jogador, permitindo uma análise mais granular do desempenho das equipes.
        Todos os dados e códigos estão disponíveis publicamente para garantir a reprodutibilidade.

        \vspace{\onelineskip}

        \noindent
        \textbf{Palavras-chave}: futebol brasileiro, inferência bayesiana, modelo Bradley-Terry, modelo Poisson, Calibração Baseada em Simulação, modelagem em nível de jogador.
    \end{otherlanguage*}
\end{resumo}
